% TOP_PROBLEM: {"id": 8000012, "title": "Perturbation theory for infinite dimensional integrable systems", "category_id": 16, "category": {"id": 16, "name": "physics", "display_name": "Mathematical Physics", "description": "Problems at the intersection of mathematics and physics.", "slug": "physics", "order_index": 16, "created_at": "2026-07-31T15:26:25.671Z"}, "status": "partially_solved", "statusQualification": "The catalogue extract has incomplete sentences and a cut-off ending. The target is reconstructed from the complete primary Problem 12, retaining the focusing multi-soliton, Toda, sine-Gordon and periodic questions. The source’s defocusing progress and periodic KAM discussion are historical context, not fresh verification of integrability or settlement."}
% FIRST_POSTED: 2026-10-11
% ENTRY_KIND: statement_only
% SUBMISSION: {"kind": "statement_only", "category": "top_problems", "problem_id": "8000012", "proposal_type": "established", "source": "Percy Deift, Some open problems in random matrix theory and the theory of integrable systems (2007), arXiv:0712.0849, Problem 12. https://arxiv.org/abs/0712.0849", "status_evidence": "Inherited source-catalogue status partially_solved preserved. Primary source accessed 11 October 2026; compared against the supplied catalogue excerpt. Historical wording and partial-result remarks were checked for scope; no comprehensive current-literature or proof-verification audit was performed.", "status_checked_on": "2026-10-11", "duplicate_check": "Full 15,773-record source registry checked for identical normalized statements; available site registry checked for identical normalized title/statement and canonical-ID collision. Canonical ID is absent from the site registry. Related canonical records 8100007. No mathematical-equivalence literature audit performed.", "attribution": "Percy Deift. Statement preparation: OpenAI GPT-6 Codex assistant, generation 11 October 2026; source-matching and notation editing by the assistant, no independent human mathematical verification. Inherited rights/provenance note: Public arXiv source TeX; arXiv reuse license NEEDS_REVIEW"}
% SOURCE_STATEMENT: "The bijective mapping properties of the scattering transform for a variety of integrable systems on the line between appropriate weighted Sobolev spaces have now been established (X. This has made it possible, in particular, to analyze (Deift-Zhou) the long-time behavior of the solution of the Cauchy problem for a variety of integrable systems in a fixed space without a ``loss of derivatives''. This in turn has made it possible to analyze perturbations of integrable systems. For example, Deift-Zhou analyzed the perturbed defocusing NLS equation iu_t+u_{xx}-2|u|^2u-\\varepsilon V(|u|)u=0 $$ u(x,t=0)=u_0(x)\\in H^{1,1}=\\{f\\in L^2\\,|\\, f', xf\\in L^2(\\mathbb{R})\\}\\,, $$ where $V(|u|)\\sim |u|^p$ as $|u|\\to0$, for some $p>2$ sufficiently large. In the perturbation theory of the linear Schr\\\"odinger equation -- when the term $-2|u|^2u$ is absent from (E:7) -- the key role in the analysis is played by the Fourier transform which diagonalizes the linear part of the equation. Such solutions do not decay uniformly in time (i.e. $\\sup_{t,x\\in\\mathbb{R}}|u(x,t)|>0$) and this complicates the analysis of the perturbed equation enormously, as the perturbation term is no longer small with respect to $2|u|^2u$ (see above). The solution of (E:7) in the focusing case in the neighborhood of a $k$-soliton for (cubic) NLS ($k\\geq2$), together with a detailed description of the long-time asymptotics, would be regarded as a very significant development in the theory of PDEs/Mathematical Physics. Here KAM methods apply and the authors show that certain finite dimensional tori corresponding to finite gap solutions in the integrable case survive under perturbation. The periodic problem is more complicated than the problem on the line because of the action of"
% Preparation and final self-audit: CONTRIBUTING.md prompts C and D; source/audit artifacts retained in the local physics_statements_2026_10_11 batch.
% Imported from: physics_statements.tex; UnsolvedMath ID 8000012
% !TeX program = lualatex
\documentclass[11pt]{article}
\usepackage[margin=25mm]{geometry}
\usepackage{fontspec}
\setmainfont{Latin Modern Roman}
\usepackage{amsmath,amssymb,mathtools}
\usepackage{unicode-math}
\setmathfont{Latin Modern Math}
\usepackage{xurl,hyperref}
\hypersetup{colorlinks=true,urlcolor=blue}
\setcounter{secnumdepth}{0}
\setlength{\parindent}{0pt}
\setlength{\parskip}{5pt}
\setlength{\emergencystretch}{3em}
\newcommand{\sref}[2]{\hyperlink{source:#1:#2}{[S#2]}}
\newcommand{\cataloguescope}[1]{\paragraph{Recorded target.}#1}
\newcommand{\R}{\mathbb R}
\newcommand{\C}{\mathbb C}
\newcommand{\N}{\mathbb N}
\newcommand{\Z}{\mathbb Z}
\newcommand{\Q}{\mathbb Q}
\newcommand{\F}{\mathbb F}
\newcommand{\eps}{\varepsilon}
\begin{document}
% UnsolvedMath ID 8000012; source catalogue code AMR-079-0012
\setcounter{section}{8000011}
\section{Perturbation theory for infinite dimensional integrable systems}\label{8000012}
{\small UnsolvedMath ID 8000012 · Mathematical Physics}
\subsection{Definitions and mathematical statement}

\paragraph{Shared notation.}$\mathbb N=\{1,2,\ldots\}$; the standard number fields have their usual meanings. All additional conventions are stated locally.


\paragraph{Definitions and notation.}An integrable evolution admits a special algebraic or spectral description allowing exact solution; a perturbation changes its equation by a small parameter-dependent term. For nonlinear Schrödinger equations, scattering coordinates may be used to compare perturbed and unperturbed dynamics. Long-time persistence, soliton behavior and asymptotic stability are different mathematical targets.

\paragraph{Catalogue wording retained for traceability.}
Develop perturbation theory for infinite-dimensional integrable systems. For the NLS example retain initial data $u_0\in H^{1,1}=\{f\in L^2(\mathbb R):f',xf\in L^2(\mathbb R)\}$ and a perturbing function $V(r)\sim r^p$ as $r\to0$, with sufficiently large $p>2$ and small positive $\varepsilon$. Beyond the source's defocusing context $iu_t+u_{xx}-2|u|^2u-\varepsilon V(|u|)u=0$, analyze the focusing sign, particularly near a $k$-soliton with $k\ge2$, and obtain detailed long-time asymptotics. Investigate persistence of Toda shock and rarefaction phenomena when exponential forces are perturbed, and sine-Gordon perturbations motivated by the Fermi--Pasta--Ulam problem. Investigate related spatially periodic Hamiltonian perturbations and surviving finite-gap tori, where whole-line dispersive decay cannot be assumed.
\paragraph{Scope and source qualification.}The catalogue extract has incomplete sentences and a cut-off ending. The target is reconstructed from the complete primary Problem 12, retaining the focusing multi-soliton, Toda, sine-Gordon and periodic questions. The source’s defocusing progress and periodic KAM discussion are historical context, not fresh verification of integrability or settlement.
\subsection{Short English statement}
Develop perturbation theory for infinite-dimensional integrable systems, including focusing NLS and multi-soliton behavior.
\subsection{Sources}
\begin{itemize}
\item[\hypertarget{source:8000012:1}{S1}] ULAM.AI and the UnsolvedMath Contributors, UnsolvedMath, record 8000012 (AMR-079-0012). Catalogue curation under CC BY 4.0; definitions and English exposition added locally.
\url{https://huggingface.co/datasets/ulamai/UnsolvedMath}.
\item[\hypertarget{source:8000012:2}{S2}] Percy Deift, Some open problems in random matrix theory and the theory of integrable systems (2007), arXiv:0712.0849, Problem 12.
\url{https://arxiv.org/abs/0712.0849}.
\end{itemize}
\label{end:8000012}
\end{document}
